\chapter{D'Artagnan Continues his Investigations}

\lettrine{A}{t} daybreak D'Artagnan saddled Furet, who had fared sumptuously all night, devouring the remainder of the oats and hay left by his companions. The musketeer sifted all he possibly could out of the host, who he found cunning, mistrustful, and devoted, body and soul, to M.~Fouquet. In order not to awaken the suspicions of this man, he carried on his fable of being a probable purchaser of some salt-mines. To have embarked for Belle-Isle at Roche-Bernard, would have been to expose himself still further to comments which had, perhaps, been already made, and would be carried to the castle. Moreover, it was singular that this traveller and his lackey should have remained a mystery to D'Artagnan, in spite of all the questions addressed by him to the host, who appeared to know him perfectly well. The musketeer then made some inquiries concerning the salt-mines, and took the road to the marshes, leaving the sea on his right, and penetrating into that vast and desolate plain which resembles a sea of mud, of which, here and there, a few crests of salt silver the undulations. Furet walked admirably, with his little nervous legs, along the foot-wide causeways which separate the salt-mines. D'Artagnan, aware of the consequences of a fall, which would result in a cold bath, allowed him to go as he liked, contenting himself with looking at, on the horizon, three rocks, that rose up like lance-blades from the bosom of the plain, destitute of verdure. Piriac, the bourgs of Batz and Le Croisic, exactly resembling each other, attracted and suspended his attention. If the traveller turned round, the better to make his observations, he saw on the other side an horizon of three other steeples, Guèrande, Le Pouliguen, and Saint-Joachim, which, in their circumference, represented a set of skittles, of which he and Furet were but the wandering ball. Piriac was the first little port on his right. He went thither, with the names of the principal salters on his lips. At the moment he reached the little port of Piriac, five large barges, laden with stone, were leaving it. It appeared strange to D'Artagnan, that stones should be leaving a country where none are found. He had recourse to all the amenity of M.~Agnan to learn from the people of the port the cause of this singular arrangement. An old fisherman replied to M.~Agnan, that the stones very certainly did not come from Piriac or the marshes.

<Where do they come from, then?> asked the musketeer.

<Monsieur, they come from Nantes and Paimbœuf.>

<Where are they going, then?>

<Monsieur, to Belle-Isle.>

<Ah! ah!> said D'Artagnan, in the same tone he had assumed to tell the printer that his character interested him; <are they building at Belle-Isle, then?>

<Why, yes, monsieur, M.~Fouquet has the walls of the castle repaired every year.>

<It is in ruins, then?>

<It is old.>

<Thank you.>

<The fact is,> said D'Artagnan to himself, <nothing is more natural; every proprietor has a right to repair his own property. It would be like telling me I was fortifying the Image-de-Notre-Dame, when I was simply obliged to make repairs. In good truth, I believe false reports have been made to his majesty, and he is very likely to be in the wrong.>

<You must confess,> continued he then, aloud, and addressing the fisherman—for his part of a suspicious man was imposed upon him by the object even of his mission—<you must confess, my dear monsieur, that these stones travel in a very curious fashion.>

<How so?> said the fisherman.

<They come from Nantes or Paimbœuf by the Loire, do they not?>

<With the tide.>

<That is convenient,—I don't say it is not; but why do they not go straight from Saint-Nazaire to Belle-Isle?>

<Eh! because the chalands (barges) are fresh-water boats, and take the sea badly,> replied the fisherman.

<That is not sufficient reason.>

<Pardon me, monsieur, one may see that you have never been a sailor,> added the fisherman, not without a sort of disdain.

<Explain to me, if you please, my good man. It appears to me that to come from Paimbœuf to Piriac, and go from Piriac to Belle-Isle, is as if we went from Roche-Bernard to Nantes, and from Nantes to Piriac.>

<By water that would be the nearest way,> replied the fisherman imperturbably.

<But there is an elbow?>

The fisherman shook his head.

<The shortest road from one place to another is a straight line,> continued D'Artagnan.

<You forget the tide, monsieur.>

<Well! take the tide.>

<And the wind.>

<Well, and the wind.>

<Without doubt; the current of the Loire carries barks almost as far as Croisic. If they want to lie by a little, or to refresh the crew, they come to Piriac along the coast; from Piriac they find another inverse current, which carries them to the Isle-Dumal, two leagues and a half.>

<Granted.>

<There the current of the Vilaine throws them upon another isle, the Isle of Hœdic.>

<I agree with that.>

<Well, monsieur, from that isle to Belle-Isle the way is quite straight. The sea, broken both above and below, passes like a canal—like a mirror between the two isles; the chalands glide along upon it like ducks upon the Loire; that's how it is.>

<It does not signify,> said the obstinate M.~Agnan; <it is a long way round.>

<Ah! yes; but M.~Fouquet will have it so,> replied, as conclusive, the fisherman, taking off his woollen cap at the enunciation of that respected name.

A look from D'Artagnan, a look as keen and piercing as a sword-blade, found nothing in the heart of the old man but a simple confidence—on his features, nothing but satisfaction and indifference. He said, <M.~Fouquet will have it so,> as he would have said, <God has willed it.>

D'Artagnan had already advanced too far in this direction; besides, the chalands being gone, there remained nothing at Piriac but a single bark—that of the old man, and it did not look fit for sea without great preparation. D'Artagnan therefore patted Furet, who, as a new proof of his charming character, resumed his march with his feet in the salt-mines, and his nose to the dry wind, which bends the furze and the broom of this country. They reached Le Croisic about five o'clock.

If D'Artagnan had been a poet, it was a beautiful spectacle: the immense strand of a league or more, the sea covers at high tide, and which, at the reflux, appears gray and desolate, strewed with polypi and seaweed, with pebbles sparse and white, like bones in some vast old cemetery. But the soldier, the politician, and the ambitious man, had no longer the sweet consolation of looking towards heaven to read there a hope or a warning. A red sky signifies nothing to such people but wind and disturbance. White and fleecy clouds upon the azure only say that the sea will be smooth and peaceful. D'Artagnan found the sky blue, the breeze embalmed with saline perfumes, and he said: <I will embark with the first tide, if it be but in a nutshell.>

At Le Croisic as at Piriac, he had remarked enormous heaps of stone lying along the shore. These gigantic walls, diminished every tide by the barges for Belle-Isle, were, in the eyes of the musketeer, the consequence and the proof of what he had well divined at Piriac. Was it a wall that M.~Fouquet was constructing? Was it a fortification that he was erecting? To ascertain that, he must make fuller observations. D'Artagnan put Furet into a stable; supped, went to bed, and on the morrow took a walk upon the port or rather upon the shingle. Le Croisic has a port of fifty feet; it has a look-out which resembles an enormous brioche (a kind of cake) elevated on a dish. The flat strand is the dish. Hundreds of barrowsful of earth amalgamated with pebbles, and rounded into cones, with sinuous passages between, are look-outs and brioches at the same time. It is so now, and it was so two hundred years ago, only the brioche was not so large, and probably there were to be seen to trellises of lath around the brioche, which constitute an ornament, planted like gardes-fous along the passages that wind towards the little terrace. Upon the shingle lounged three or four fishermen talking about sardines and shrimps. D'Artagnan, with his eyes animated by a rough gayety, and a smile upon his lips, approached these fishermen.

<Any fishing going on to-day?> said he.

<Yes, monsieur,> replied one of them, <we are only waiting for the tide.>

<Where do you fish, my friends?>

<Upon the coasts, monsieur.>

<Which are the best coasts?>

<Ah, that is all according. The tour of the isles, for example?>

<Yes, but they are a long way off, those isles, are they not?>

<Not very; four leagues.>

<Four leagues! That is a voyage.>

The fishermen laughed in M.~Agnan's face.

<Hear me, then,> said the latter with an air of simple stupidity; <four leagues off you lose sight of land, do you not?>

<Why, not always.>

<Ah, it is a long way—too long, or else I would have asked you to take me aboard, and to show me what I have never seen.>

<What is that?>

<A live sea-fish.>

<Monsieur comes from the province?> said a fisherman.

<Yes, I come from Paris.>

The Breton shrugged his shoulders; then:

<Have you ever seen M.~Fouquet in Paris?> asked he.

<Often,> replied D'Artagnan.

<Often!> repeated the fishermen, closing their circle round the Parisian. <Do you know him?>

<A little; he is the intimate friend of my master.>

<Ah!> said the fishermen, in astonishment.

<And,> said D'Artagnan, <I have seen all his châteaux of Saint Mande, of Vaux, and his hotel in Paris.>

<Is that a fine place?>

<Superb.>

<It is not so fine a place as Belle-Isle,> said the fisherman.

<Bah!> cried M.~d'Artagnan, breaking into a laugh so loud that he angered all his auditors.

<It is very plain that you have never seen Belle-Isle,> said the most curious of the fishermen. <Do you know that there are six leagues of it, and that there are such trees on it as cannot be equalled even at Nates-sur-le-Fosse?>

<Trees in the sea!> cried D'Artagnan; <well, I should like to see them.>

<That can be easily done; we are fishing at the Isle de Hœdic—come with us. From that place you will see, as a Paradise, the black trees of Belle-Isle against the sky; you will see the white line of the castle, which cuts the horizon of the sea like a blade.>

<Oh,> said D'Artagnan, <that must be very beautiful. But do you know there are a hundred belfries at M.~Fouquet's château of Vaux?>

The Breton raised his head in profound admiration, but he was not convinced. <A hundred belfries! Ah, that may be; but Belle-Isle is finer than that. Should you like to see Belle-Isle?>

<Is that possible?> asked D'Artagnan.

<Yes, with permission of the governor.>

<But I do not know the governor.>

<As you know M.~Fouquet, you can tell your name.>

<Oh, my friends, I am not a gentleman.>

<Everybody enters Belle-Isle,> continued the fisherman in his strong, pure language, <provided he means no harm to Belle-Isle or its master.>

A slight shudder crept over the body of the musketeer. <That is true,> thought he. Then recovering himself, <If I were sure,> said he, <not to be sea-sick.>

<What, upon her?> said the fisherman, pointing with pride to his pretty round-bottomed bark.

<Well, you almost persuade me,> cried M.~Agnan; <I will go and see Belle-Isle, but they will not admit me.>

<We shall enter, safe enough.>

<You! What for?>

<Why, dame! to sell fish to the corsairs.>

<Ha! Corsairs—what do you mean?>

<Well, I mean that M.~Fouquet is having two corsairs built to chase the Dutch and the English, and we sell our fish to the crews of those little vessels.>

<Come, come!> said D'Artagnan to himself—<better and better. A printing-press, bastions, and corsairs! Well, M.~Fouquet is not an enemy to be despised, as I presumed to fancy. He is worth the trouble of travelling to see him nearer.>

<We set out at half-past five,> said the fisherman gravely.

<I am quite ready, and I will not leave you now.> So D'Artagnan saw the fishermen haul their barks to meet the tide with a windlass. The sea rose; M.~Agnan allowed himself to be hoisted on board, not without sporting a little fear and awkwardness, to the amusement of the young beach-urchins who watched him with their large intelligent eyes. He laid himself down upon a folded sail, not interfering with anything whilst the bark prepared for sea; and, with its large square sail, it was fairly out within two hours. The fishermen, who prosecuted their occupation as they proceeded, did not perceive that their passenger had not become pale, neither groaned nor suffered; that in spite of that horrible tossing and rolling of the bark, to which no hand imparted direction, the novice passenger had preserved his presence of mind and his appetite. They fished, and their fishing was sufficiently fortunate. To lines bated with prawn, soles came, with numerous gambols, to bite. Two nets had already been broken by the immense weight of congers and haddocks; three sea-eels ploughed the hold with their slimy folds and their dying contortions. D'Artagnan brought them good luck; they told him so. The soldier found the occupation so pleasant, that he put his hand to the work—that is to say, to the lines—and uttered roars of joy, and mordioux enough to have astonished his musketeers themselves, every time that a shock given to his line by the captured fish required the play of the muscles of his arm, and the employment of his best dexterity. The party of pleasure had made him forget his diplomatic mission. He was struggling with a very large conger, and holding fast with one hand to the side of the vessel, in order to seize with the other the gaping jowl of his antagonist, when the master said to him, <Take care they don't see you from Belle-Isle!>

These words produced the same effect upon D'Artagnan as the hissing of the first bullet on a day of battle; he let go of both line and conger, which, dragging each other, returned again to the water. D'Artagnan perceived, within half a league at most, the blue and marked profile of the rocks of Belle-Isle, dominated by the majestic whiteness of the castle. In the distance, the land with its forests and verdant plains; cattle on the grass. This was what first attracted the attention of the musketeer. The sun darted its rays of gold upon the sea, raising a shining mist round this enchanted isle. Little could be seen of it, owing to this dazzling light, but the salient points; every shadow was strongly marked, and cut with bands of darkness the luminous fields and walls. <Eh! eh!> said D'Artagnan, at the aspect of those masses of black rocks, <these are fortifications which do not stand in need of any engineer to render a landing difficult. How the devil can a landing be effected on that isle which God has defended so completely?>

<This way,> replied the patron of the bark, changing the sail, and impressing upon the rudder a twist which turned the boat in the direction of a pretty little port, quite coquettish, round, and newly battlemented.

<What the devil do I see yonder?> said D'Artagnan.

<You see Locmaria,> replied the fisherman.

<Well, but there?>

<That is Bangor.>

<And further on?>

<Sauzon, and then Le Palais.>

<\swear{Mordioux!} It is a world. Ah! there are some soldiers.>

<There are seventeen hundred men in Belle-Isle, monsieur,> replied the fisherman, proudly. <Do you know that the least garrison is of twenty companies of infantry?>

<\swear{Mordioux!}> cried D'Artagnan, stamping with his foot. <His majesty was right enough.>

They landed. 